\paragraph{Uniform finite first-order comparison.}
The characteristic hypothesis can be checked without rate-dependent
hidden constants. For the parameters of
Theorem~\ref{thm:projective-quadratic-line}, take the derivative-weighted
support of~\cite[Eq.~(60)]{dkt2026} with
\[
 m=4,\qquad B_\partial=2,\qquad B=B_{\rm jet}=2A_0,
 \qquad H=B_Z=116B.
\]
Its coefficient count is $G=3B^2$ and its local-rank bound is
$R=3+6+8+6=23$. Direct expansion gives
\[
 5G-116N=8(p-5)^4+168(p-5)^3+1148(p-5)^2+2648(p-5)+308>0.
\]
Thus $(H-B+1)G>(H+1)23N$, which suffices for the graded challenge
test of~\cite[Proposition~5.10]{dkt2026}. Reconstruction requires only
$p>\max\{2,B_\partial\}$.

For clarity, the resulting bounds are uniform over every received word
and every affine received line on this evaluation set. In the notation
of~\cite[Lemma~5.6]{dkt2026}, one has
\[
 \lambda\le3N/B,\quad F_{\rm reg}\le5B,\quad
 S\le3B,\quad H_s=348B,\quad J_1\le1866B^2.
\]
Choose both incidence thresholds $L=L_0=\lfloor A_0/2\rfloor$.
The remaining incidence ratios are at most $4N/B$ and $5N/B$.
Using $B<N$, $B^2<8N$, and $\Psi_2(S)\le24B$, the list budget is
at most $18N$, and the full-agreement-set MCA exceptional budget is
at most
\[
 22392N^2+25N^2+16716N+33411NB<75000N^2.
\]
The lower bound $M>\tfrac14N^{3/2}$ therefore lies in the same
explicit large-characteristic first-order setting, though it does not
match this quadratic upper bound.
